\documentclass[11pt]{article}
\usepackage[a4paper,margin=22mm]{geometry}
\usepackage{fontspec}
\setmainfont{DejaVu Serif}
\setsansfont{DejaVu Sans}
\usepackage{microtype}
\usepackage{xcolor}
\usepackage{hyperref}
\usepackage{enumitem}
\usepackage{parskip}
\definecolor{Accent}{HTML}{971D32}
\definecolor{Muted}{HTML}{666666}
\hypersetup{colorlinks=true,urlcolor=Accent,linkcolor=Accent}
\setlist{leftmargin=1.6em,itemsep=0.3em,topsep=0.35em}
\setcounter{tocdepth}{2}
\renewcommand{\contentsname}{Contents}
\newcommand{\sectionrule}{\par\vspace{-0.5em}\noindent\textcolor{Accent}{\rule{\linewidth}{0.6pt}}\par}
\begin{document}

\begin{center}
{\sffamily\bfseries\color{Accent} TOEFL WORD OF THE DAY\par}
\vspace{8mm}
{\fontsize{42}{48}\selectfont\bfseries impose\par}
\vspace{2mm}
{\Large\sffamily\color{Accent} /ɪmˈpoʊz/\par}
\vspace{2mm}
{\small\sffamily Local audio phrase: impose\par}
{\small\sffamily Audio file: audios/impose.mp3\par}
\vspace{2mm}
{\large\sffamily\color{Muted} transitive verb\par}
\vspace{5mm}
{\sffamily officially introduce and enforce \textbullet\ force something unwanted on someone \textbullet\ burden someone with your presence or demands\par}
\vspace{6mm}
{\large Impose usually involves making people accept something they did not freely choose, such as a rule, tax, burden, belief, or demand.\par}
\vspace{6mm}
\begin{minipage}{0.84\textwidth}
\itshape “The government imposed stricter limits on industrial pollution.”
\end{minipage}\par
\vspace{10mm}
\textbf{Date:} August 26th 2026\quad\textbullet\quad \textbf{Level:} B1--mid-B2 TOEFL
\vfill
Created and compiled by \textbf{Amir Kooshky}\par
\vspace{2mm}
\href{https://t.me/KooshkyTOEFL}{Telegram Channel}\quad\textbullet\quad
\href{https://www.instagram.com/kooshkytoefl}{Instagram}\quad\textbullet\quad
\href{https://t.me/KooshkyTOEFL\_pv}{Personal Telegram}
\end{center}
\newpage
\tableofcontents
\newpage

\section{Meanings and usage}
\sectionrule
\subsection{Meaning 1: Officially introduce and enforce}
\textbf{Part of speech:} transitive verb

\textbf{IPA:} /ɪmˈpoʊz/

\textbf{Pronunciation phrase:} impose

\textbf{Local audio file:} audios/impose.mp3

\textbf{Register:} neutral to formal; very common in academic, legal, political, economic, and policy contexts

\textbf{Definition:} to officially introduce a rule, tax, punishment, restriction, or other requirement and make people obey or accept it

\textbf{In simpler words:} officially make a rule or requirement apply

\textbf{Typical pattern:} impose + tax, fine, ban, restriction, limit, penalty, sentence, or requirement

\subsubsection{Natural examples}
\begin{enumerate}
  \item The government imposed stricter limits on industrial pollution.
  \item The court imposed a substantial fine on the company.
  \item Several countries imposed new taxes on imported goods.
  \item The university imposed a temporary ban on parking near the construction site.
  \item Authorities imposed restrictions on water use during the drought.
  \item The treaty imposed strict limits on the country's military activities.
  \item The judge imposed a sentence of three years in prison.
  \item New regulations impose additional safety requirements on manufacturers.
\end{enumerate}

\subsubsection{Nuance and implication}
\textbf{Central nuance:} Impose emphasizes authority and lack of free choice: a government, court, institution, or other authority makes something apply to people.

\textbf{Usual implication:} The people affected may not agree with or welcome the rule, tax, restriction, or punishment.

\textbf{Compared with recommend:} To recommend something is to advise people to do it. To impose it is to make it compulsory or officially apply it.

\subsubsection{Grammar notes}
\textbf{How to use it:} Put the rule, tax, punishment, or restriction directly after impose.

\textbf{What can follow it:} Common objects include tax, fine, penalty, ban, restriction, limit, requirement, condition, sanction, and sentence.

\textbf{Common preposition:} When naming the people or organization affected, use on or upon: impose a tax on consumers, impose restrictions on factories.

\textbf{Position in a sentence:} The passive form is also common: A fine was imposed on the company.

\subsubsection{Useful patterns}
\textbf{Pattern:} impose + tax or fee + on + person or group

\textbf{Use:} Use this when an authority officially requires people to pay money.

\textbf{Example:} The city imposed a new fee on property developers.

\textbf{Pattern:} impose + restriction or limit + on + activity

\textbf{Use:} Use this when an authority officially limits what people can do.

\textbf{Example:} The government imposed restrictions on commercial fishing.

\textbf{Pattern:} impose + fine or penalty + on + person or organization

\textbf{Use:} Use this when an authority officially punishes someone.

\textbf{Example:} Regulators imposed a heavy fine on the company.

\textbf{Pattern:} impose + ban

\textbf{Use:} Use this when an authority officially prohibits something.

\textbf{Example:} The city imposed a temporary ban on outdoor fires.

\textbf{Pattern:} be imposed on + person, group, or activity

\textbf{Use:} Use the passive form when the rule or punishment is the main focus.

\textbf{Example:} Strict limits were imposed on groundwater use.

\subsubsection{Common wrong usage}
\paragraph{Correction 1}
\textbf{Wrong:} The government imposed to companies stricter regulations.

\textbf{Correct:} The government imposed stricter regulations on companies.

\textbf{Why:} Put the thing being imposed directly after impose, and use on before the people or organizations affected.

\paragraph{Correction 2}
\textbf{Wrong:} The court imposed the company a fine.

\textbf{Correct:} The court imposed a fine on the company.

\textbf{Why:} The normal pattern is impose something on someone.

\paragraph{Correction 3}
\textbf{Wrong:} The government imposed people to follow the new rule.

\textbf{Correct:} The government required people to follow the new rule.

\textbf{Why:} Do not use impose + person + to + verb. Use impose a rule on someone, or use require someone to do something.

\subsection{Meaning 2: Force something unwanted on someone}
\textbf{Part of speech:} transitive verb

\textbf{IPA:} /ɪmˈpoʊz/

\textbf{Pronunciation phrase:} impose

\textbf{Local audio file:} audios/impose.mp3

\textbf{Register:} neutral to formal; common in academic discussions of policy, economics, society, education, and ethics

\textbf{Definition:} to force an unwanted burden, idea, condition, or responsibility on someone

\textbf{In simpler words:} force someone to accept or deal with something

\textbf{Typical pattern:} impose + burden, cost, pressure, condition, belief, responsibility, or demand + on + someone

\subsubsection{Natural examples}
\begin{enumerate}
  \item Parents should be careful not to impose their own career ambitions on their children.
  \item The new reporting requirements impose a significant burden on small businesses.
  \item Critics argued that the policy imposed unnecessary costs on local communities.
  \item A teacher should not impose personal political beliefs on students.
  \item The project imposed additional responsibilities on employees who were already overworked.
  \item Rapid population growth can impose severe pressure on public services.
  \item The agreement imposed several difficult conditions on the smaller company.
  \item The disaster imposed enormous financial costs on the region.
\end{enumerate}

\subsubsection{Nuance and implication}
\textbf{Central nuance:} This sense focuses on making someone carry, accept, or deal with something they did not freely choose.

\textbf{Usual implication:} The imposed thing is often difficult, unwanted, restrictive, expensive, or unfair.

\textbf{Compared with offer:} Something offered can be accepted or refused. Something imposed is forced on the person or group.

\subsubsection{Grammar notes}
\textbf{How to use it:} Put the burden, idea, condition, or responsibility directly after impose.

\textbf{What can follow it:} Common objects include burden, cost, pressure, condition, responsibility, demand, belief, value, and expectation.

\textbf{Common preposition:} Use on or upon before the person or group receiving the unwanted thing: impose a burden on taxpayers.

\textbf{Position in a sentence:} The passive is common when focusing on the affected person or group: Extra responsibilities were imposed on staff.

\subsubsection{Useful patterns}
\textbf{Pattern:} impose a burden on + person or group

\textbf{Use:} Use this when something creates extra difficulty, work, or expense.

\textbf{Example:} The paperwork imposes a heavy burden on small businesses.

\textbf{Pattern:} impose costs on + person or group

\textbf{Use:} Use this when a decision causes people to pay or lose money.

\textbf{Example:} Environmental damage can impose substantial costs on future generations.

\textbf{Pattern:} impose pressure on + person, system, or resource

\textbf{Use:} Use this when something creates strain or difficulty.

\textbf{Example:} Rapid urban growth imposes pressure on transportation systems.

\textbf{Pattern:} impose beliefs or values on + someone

\textbf{Use:} Use this when someone tries to force another person to accept particular ideas.

\textbf{Example:} Schools should educate students without imposing a particular political ideology on them.

\textbf{Pattern:} impose conditions on + someone

\textbf{Use:} Use this when someone must accept certain requirements.

\textbf{Example:} The lender imposed strict conditions on the borrower.

\subsubsection{Common wrong usage}
\paragraph{Correction 1}
\textbf{Wrong:} Parents should not impose their children their beliefs.

\textbf{Correct:} Parents should not impose their beliefs on their children.

\textbf{Why:} Use impose something on someone.

\paragraph{Correction 2}
\textbf{Wrong:} The policy imposed many pressures to small businesses.

\textbf{Correct:} The policy imposed considerable pressure on small businesses.

\textbf{Why:} Use on, not to. Pressure is also commonly uncountable in this meaning.

\paragraph{Correction 3}
\textbf{Wrong:} She imposed me to accept her opinion.

\textbf{Correct:} She tried to impose her opinion on me.

\textbf{Why:} Do not use impose someone to do something. Use impose something on someone.

\subsection{Meaning 3: Burden someone with your presence or demands}
\textbf{Part of speech:} transitive verb

\textbf{IPA:} /ɪmˈpoʊz/

\textbf{Pronunciation phrase:} impose

\textbf{Local audio file:} audios/impose.mp3

\textbf{Register:} neutral; often used politely when discussing favors, hospitality, personal space, or unwanted attention

\textbf{Definition:} to force your presence, attention, or demands on someone when they may not want them

\textbf{In simpler words:} cause someone inconvenience by asking too much or staying when not wanted

\textbf{Typical pattern:} impose on + someone, or impose yourself on + someone

\subsubsection{Natural examples}
\begin{enumerate}
  \item I did not want to impose on my neighbors by asking them to drive me to the airport.
  \item She felt that staying for another week would impose on her hosts.
  \item He was careful not to impose himself on colleagues who were already extremely busy.
  \item I hope I am not imposing, but could I use your phone for a moment?
  \item They offered us dinner, but we did not want to impose on their hospitality.
  \item She disliked imposing on friends for help unless it was absolutely necessary.
  \item The visitor imposed himself on the family and stayed far longer than expected.
\end{enumerate}

\subsubsection{Nuance and implication}
\textbf{Central nuance:} This sense often expresses concern that your request, presence, or behavior may inconvenience another person.

\textbf{Usual implication:} It is frequently used in polite language when someone wants to avoid taking advantage of another person's kindness.

\textbf{Compared with respect someone's space:} Respecting someone's space means avoiding unwanted demands or intrusion. Imposing on someone means creating an unwanted burden through your presence or requests.

\subsubsection{Grammar notes}
\textbf{How to use it:} Use impose on someone when your request or presence may burden them. Use impose yourself on someone when you force your company or attention on them.

\textbf{What can follow it:} After impose on, name the person or group affected. You can also say impose on someone's hospitality or generosity.

\textbf{Common preposition:} On is essential in this pattern: impose on someone.

\textbf{Position in a sentence:} This meaning often appears in negative or polite expressions such as I don't want to impose or I hope I'm not imposing.

\subsubsection{Useful patterns}
\textbf{Pattern:} impose on + someone

\textbf{Use:} Use this when asking a favor or staying with someone might inconvenience them.

\textbf{Example:} I do not want to impose on you, so I can find a hotel instead.

\textbf{Pattern:} impose on someone's hospitality

\textbf{Use:} Use this when you worry that you are accepting too much kindness from a host.

\textbf{Example:} We decided to leave early rather than impose on their hospitality.

\textbf{Pattern:} impose yourself on + someone

\textbf{Use:} Use this when a person forces their presence or attention on someone else.

\textbf{Example:} He had a habit of imposing himself on people who wanted to be left alone.

\textbf{Pattern:} I hope I'm not imposing

\textbf{Use:} Use this polite expression before or after making a request that may inconvenience someone.

\textbf{Example:} I hope I'm not imposing, but could I stay here tonight?

\subsubsection{Common wrong usage}
\paragraph{Correction 1}
\textbf{Wrong:} I don't want to impose you.

\textbf{Correct:} I don't want to impose on you.

\textbf{Why:} When impose means inconvenience someone through a request or your presence, use impose on someone.

\paragraph{Correction 2}
\textbf{Wrong:} He imposed himself to his neighbors.

\textbf{Correct:} He imposed himself on his neighbors.

\textbf{Why:} The correct pattern is impose yourself on someone.

\paragraph{Correction 3}
\textbf{Wrong:} I hope I don't impose to ask for another favor.

\textbf{Correct:} I hope I'm not imposing by asking for another favor.

\textbf{Why:} A natural pattern is be imposing by + -ing form.

\section{Word family}
\sectionrule
\subsection{imposition}
\textbf{Part of speech:} countable noun

\textbf{IPA:} /ˌɪmpəˈzɪʃən/

\textbf{Definition:} an unfair or inconvenient demand or burden placed on someone

\textbf{Relationship to the headword:} This noun can refer to a burden or demand that has been imposed on someone.

\textbf{Grammar pattern:} an imposition on + person or group

\textbf{Usage note:} It is especially common when talking about inconvenience or an unreasonable demand.

\subsubsection{Examples}
\begin{enumerate}
  \item Requiring employees to work every weekend would be a considerable imposition.
  \item She apologized for the imposition and thanked her neighbors for helping.
  \item The additional paperwork was seen as an unnecessary imposition.
\end{enumerate}

\subsection{imposing}
\textbf{Part of speech:} adjective

\textbf{IPA:} /ɪmˈpoʊzɪŋ/

\textbf{Definition:} large, impressive, and likely to attract attention because of appearance or size

\textbf{Relationship to the headword:} This adjective is historically related to impose but has developed the separate common meaning of impressive in size or appearance.

\textbf{Grammar pattern:} imposing + building, structure, figure, entrance, or appearance

\textbf{Usage note:} Do not assume that imposing means forcing something in this adjective use.

\subsubsection{Examples}
\begin{enumerate}
  \item The museum occupies an imposing stone building in the center of the city.
  \item The mountain presents an imposing sight from the valley below.
\end{enumerate}

\section{Common collocations}
\sectionrule
\subsection{impose a tax}
\textbf{Applies to:} Officially introduce and enforce

\textbf{Meaning:} officially require a tax to be paid

\textbf{Pattern:} impose a tax on + person, product, or activity

\textbf{Example:} The government imposed a tax on luxury goods.

\subsection{impose a fine}
\textbf{Applies to:} Officially introduce and enforce

\textbf{Meaning:} officially require someone to pay money as punishment

\textbf{Pattern:} impose a fine on + person or organization

\textbf{Example:} The regulator imposed a fine on the company for violating safety rules.

\subsection{impose a penalty}
\textbf{Applies to:} Officially introduce and enforce

\textbf{Meaning:} officially give a punishment

\textbf{Example:} The law allows courts to impose severe penalties for repeated violations.

\subsection{impose restrictions}
\textbf{Applies to:} Officially introduce and enforce

\textbf{Meaning:} officially limit what people are allowed to do

\textbf{Pattern:} impose restrictions on + activity or group

\textbf{Example:} Authorities imposed restrictions on construction near the protected area.

\subsection{impose a ban}
\textbf{Applies to:} Officially introduce and enforce

\textbf{Meaning:} officially prohibit something

\textbf{Pattern:} impose a ban on + activity or product

\textbf{Example:} The government imposed a ban on the sale of the chemical.

\subsection{impose limits}
\textbf{Applies to:} Officially introduce and enforce

\textbf{Meaning:} officially restrict the amount or extent of something

\textbf{Pattern:} impose limits on + activity or quantity

\textbf{Example:} The law imposes strict limits on industrial emissions.

\subsection{impose sanctions}
\textbf{Applies to:} Officially introduce and enforce

\textbf{Meaning:} officially introduce economic or political penalties

\textbf{Pattern:} impose sanctions on + country, organization, or individual

\textbf{Example:} Several countries imposed economic sanctions on the regime.

\subsection{impose a burden}
\textbf{Applies to:} Force something unwanted on someone

\textbf{Meaning:} make someone carry an extra difficulty, responsibility, or cost

\textbf{Pattern:} impose a burden on + person or group

\textbf{Example:} The new rules impose an additional burden on small businesses.

\subsection{impose costs}
\textbf{Applies to:} Force something unwanted on someone

\textbf{Meaning:} cause someone or society to bear financial or other costs

\textbf{Pattern:} impose costs on + person, group, or society

\textbf{Example:} Pollution imposes substantial costs on society.

\subsection{impose conditions}
\textbf{Applies to:} Force something unwanted on someone

\textbf{Meaning:} require someone to accept particular terms

\textbf{Pattern:} impose conditions on + person, agreement, or activity

\textbf{Example:} The bank imposed strict conditions on the loan.

\subsection{impose one's views}
\textbf{Applies to:} Force something unwanted on someone

\textbf{Meaning:} force other people to accept your opinions

\textbf{Pattern:} impose one's views on + someone

\textbf{Example:} Teachers should avoid imposing their personal views on students.

\subsection{impose on someone}
\textbf{Applies to:} Burden someone with your presence or demands

\textbf{Meaning:} cause someone inconvenience through your requests or presence

\textbf{Example:} I do not want to impose on you by staying another night.

\subsection{impose on someone's hospitality}
\textbf{Applies to:} Burden someone with your presence or demands

\textbf{Meaning:} take unfair advantage of someone's willingness to host or help you

\textbf{Example:} We did not want to impose on our hosts' hospitality.

\section{Synonyms and antonyms}
\sectionrule
\subsection{Close synonyms}
\subsubsection{enforce}
\textbf{Part of speech:} transitive verb

\textbf{Applies to:} Officially introduce and enforce

\textbf{Shared meaning:} Both involve authority and compulsory rules.

\textbf{Difference:} Enforce means make sure an existing law or rule is obeyed. Impose means introduce or place the rule, restriction, or punishment on people in the first place.

\textbf{Example:} Police are responsible for enforcing traffic laws.

\subsubsection{levy}
\textbf{Part of speech:} transitive verb

\textbf{Applies to:} Officially introduce and enforce

\textbf{Shared meaning:} Both can mean officially require a tax or financial charge.

\textbf{Difference:} Levy is a more specialized formal verb used mainly for officially imposing taxes, fees, or other financial charges. Impose can be used much more broadly.

\textbf{Example:} The city levied a tax on vacant properties.

\subsubsection{force}
\textbf{Part of speech:} transitive verb

\textbf{Applies to:} Force something unwanted on someone

\textbf{Shared meaning:} Both involve lack of free choice.

\textbf{Difference:} Force is broader and often focuses on making a person perform an action. Impose usually focuses on placing a rule, burden, belief, condition, or cost on someone.

\textbf{Example:} The circumstances forced the company to reduce production.

\subsubsection{burden}
\textbf{Part of speech:} transitive verb

\textbf{Applies to:} Force something unwanted on someone

\textbf{Shared meaning:} Both can describe causing someone to carry an unwanted difficulty.

\textbf{Difference:} Burden emphasizes giving someone difficult responsibilities, costs, or problems. Impose is broader and can also refer to rules, beliefs, restrictions, or conditions.

\textbf{Example:} The complicated process burdens small businesses with unnecessary paperwork.

\subsubsection{intrude}
\textbf{Part of speech:} intransitive verb

\textbf{Applies to:} Burden someone with your presence or demands

\textbf{Shared meaning:} Both can describe unwanted interference with another person.

\textbf{Difference:} Intrude emphasizes entering someone's space, privacy, or activity without permission. Impose on emphasizes causing inconvenience by asking too much or taking advantage of someone's hospitality.

\textbf{Example:} I did not want to intrude on their private conversation.

\subsection{Direct or useful antonyms}
\subsubsection{lift}
\textbf{Part of speech:} transitive verb

\textbf{Applies to:} Officially introduce and enforce

\textbf{Opposition:} To impose a restriction or ban is to introduce it; to lift it is to officially remove it.

\textbf{Usage note:} This is a strong opposite specifically for restrictions, bans, sanctions, and similar measures.

\textbf{Example:} Authorities lifted the restrictions after conditions improved.

\subsubsection{remove}
\textbf{Part of speech:} transitive verb

\textbf{Applies to:} Officially introduce and enforce

\textbf{Opposition:} Impose can mean place a requirement or restriction on someone, while remove means take that requirement or restriction away.

\textbf{Example:} The government removed several unnecessary restrictions.

\section{Words learners may confuse}
\sectionrule
\subsection{enforce}
\textbf{Part of speech:} transitive verb

\textbf{Why learners confuse them:} Both commonly appear in discussions of laws, rules, and government authority.

\textbf{Key difference:} Impose means introduce or place a rule, tax, punishment, or restriction on someone. Enforce means make sure an existing rule or law is obeyed.

\textbf{impose:} The city imposed a new parking restriction.

\textbf{enforce:} Police began enforcing the new parking restriction on Monday.

\subsection{propose}
\textbf{Part of speech:} transitive verb

\textbf{Why learners confuse them:} The words look somewhat similar but represent very different stages of decision-making.

\textbf{Key difference:} Impose means make something apply without giving the affected people a choice. Propose means suggest something for consideration.

\textbf{impose:} The government imposed a new tax.

\textbf{propose:} The government proposed a new tax, but parliament had not yet approved it.

\subsection{imposing}
\textbf{Part of speech:} verb versus adjective

\textbf{Why learners confuse them:} Imposing looks like the -ing form of impose but also has an established adjective meaning that learners may not expect.

\textbf{Key difference:} Impose is a verb meaning force or officially place something on someone. Imposing is commonly an adjective meaning large, impressive, and visually powerful.

\textbf{impose:} The court imposed a heavy fine on the company.

\textbf{imposing:} The courthouse is an imposing stone building.

\section{Etymology}
\sectionrule
\textbf{Origin:} Impose comes through French from Latin elements meaning to put or place upon.

\textbf{Development:} The idea of placing something upon another thing developed into the modern meanings of placing a rule, burden, duty, tax, or demand on someone.

\textbf{Memory link:} Think of impose as putting something on someone: a tax, rule, burden, belief, or demand.

\section{Practice exercises}
\sectionrule
\subsection{Word family choice}
Choose the best member of the word family for each blank.

\begin{enumerate}
  \item Requiring employees to attend unpaid meetings every weekend would be a serious \_\_\_\_\_.
  \item The museum is housed in an \_\_\_\_\_ nineteenth-century building.
  \item The government decided to \_\_\_\_\_ stricter limits on industrial emissions.
\end{enumerate}

\subsection{Correct or incorrect?}
Decide whether each sentence is correct. Correct the inaccurate ones.

\begin{enumerate}
  \item The court imposed a heavy fine on the company.
  \item The government imposed citizens to follow the new rule.
  \item The new regulations impose additional costs on manufacturers.
  \item Parents should not impose their ambitions to their children.
  \item I do not want to impose on you by asking for another favor.
\end{enumerate}

\subsection{Collocation practice}
Complete each sentence with a natural collocation.

\begin{enumerate}
  \item The government imposed new restrictions \_\_\_\_\_ industrial water use.
  \item The court decided to impose a substantial \_\_\_\_\_ on the company.
  \item Complex regulations can impose a heavy \_\_\_\_\_ on small businesses.
  \item Teachers should avoid imposing their personal views \_\_\_\_\_ students.
  \item I hope I am not imposing \_\_\_\_\_ you by staying another night.
\end{enumerate}

\subsection{Complete answer key}
\subsubsection{Word family choice}
\begin{enumerate}
  \item \textbf{imposition.} Requiring employees to attend unpaid meetings every weekend would be a serious imposition. \emph{Use the noun imposition for an unreasonable or inconvenient burden.}
  \item \textbf{imposing.} The museum is housed in an imposing nineteenth-century building. \emph{The adjective imposing means large and impressive in appearance.}
  \item \textbf{impose.} The government decided to impose stricter limits on industrial emissions. \emph{Use the base verb impose after to.}
\end{enumerate}

\subsubsection{Correct or incorrect?}
\begin{enumerate}
  \item \textbf{Correct} \emph{Impose something on someone is the standard pattern.}
  \item \textbf{Incorrect}. The government required citizens to follow the new rule. \emph{Do not use impose + person + to + verb.}
  \item \textbf{Correct} \emph{Impose costs on someone means cause them to bear those costs.}
  \item \textbf{Incorrect}. Parents should not impose their ambitions on their children. \emph{Use impose something on someone.}
  \item \textbf{Correct} \emph{Impose on someone means inconvenience or burden them through your requests or presence.}
\end{enumerate}

\subsubsection{Collocation practice}
\begin{enumerate}
  \item \textbf{on.} The government imposed new restrictions on industrial water use. \emph{Use impose restrictions on something.}
  \item \textbf{fine.} The court decided to impose a substantial fine on the company. \emph{Impose a fine is a common legal collocation.}
  \item \textbf{burden.} Complex regulations can impose a heavy burden on small businesses. \emph{Impose a burden on someone means cause them extra difficulty, work, or expense.}
  \item \textbf{on.} Teachers should avoid imposing their personal views on students. \emph{Use impose one's views on someone.}
  \item \textbf{on.} I hope I am not imposing on you by staying another night. \emph{Impose on someone means cause them inconvenience through your presence or requests.}
\end{enumerate}

\vfill
\begin{center}
\small\color{Muted} Created and compiled by \textbf{Amir Kooshky}\\
\href{https://t.me/KooshkyTOEFL}{Telegram Channel}\quad\textbullet\quad
\href{https://www.instagram.com/kooshkytoefl}{Instagram}\quad\textbullet\quad
\href{https://t.me/KooshkyTOEFL\_pv}{Personal Telegram}
\end{center}
\end{document}
